\chapter{\nt{GLUT} e \nt{GABA} juntos}
\label{sec:gaba}

\section{As regras do jogo}

Uma rede só de neurônios \nt{GLUT} não sabe escolher, zerar a memória nem impedir que a atividade vire avalanche, e tudo por causa da monotonicidade (proposição~\ref{prop:monotone}). Vamos acrescentar a ela o segundo tipo de neurônio mais numeroso, o \nt{GABA}. As regras são as mesmas: só o que existe num organismo vivo, e nenhum sinal de clock.

O modelo de neurônio é o mesmo de~\eqref{eq:glutneuron}, mas o disparo de um neurônio \nt{GABA} não eleva a tensão do destinatário, e sim a abaixa. Agora os pesos têm dois sinais, e o sinal é definido pelo remetente, regra~\eqref{eq:dalesign}.

Alguns parâmetros do circuito. Os neurônios \nt{GABA} são de um quinto a um quarto dos neurônios do córtex~\cite{Hendry1987}. Suas saídas são curtas: eles inibem os vizinhos, não regiões distantes. A inibição chega em um milissegundo e dura alguns milissegundos. Sem entrada, o neurônio \nt{GABA} fica calado: para inibir alguém, ele mesmo precisa ser excitado, em geral por neurônios \nt{GLUT} da mesma rede.

Por isso, uma conexão negativa de um neurônio \nt{GLUT} $A$ para um neurônio $B$ tem de passar por um intermediário: $A$ excita um neurônio \nt{GABA}, e este inibe $B$. A troca de sinal custa um neurônio e um atraso, cerca de um milissegundo.

Na inibição importa não só \emph{de quem} vem a conexão, mas também \emph{onde} ela se instala: o local de contato determina em grande parte o que a conexão faz~\cite{Markram2004}. As saídas \nt{GLUT} se instalam principalmente nos ramos de entrada do destinatário, seus \emph{dendritos}, e levam dados: cada entrada dessas é uma parcela da soma. Uma saída no corpo do neurônio age sobre a soma inteira: assim muitos neurônios \nt{GABA} rápidos dividem a resposta do destinatário e definem quando ele vai disparar. Uma saída no local onde nasce o disparo do destinatário é a última palavra, um veto à saída inteira de uma vez. E uma saída no terminal do fio de outro neurônio muda a probabilidade de liberação $p$ de uma única linha de entrada, sem tocar nas demais~\cite{Rudomin1999}; é assim, em particular, que funciona a porta da dor no capítulo~\ref{sec:pain}. Fora do cérebro, as saídas se instalam em fibras musculares (capítulo~\ref{sec:muscles}) e em órgãos (capítulo~\ref{sec:chemoutput}), e algumas não se instalam em lugar nenhum e liberam o neurotransmissor no volume do tecido ou no sangue (capítulos~\ref{sec:serocomputer} e~\ref{sec:chemoutput}).

\section{NÃO e veto}

Dá para fazer um NÃO agora? Quase. Um neurônio que trabalha com a entrada calada precisa tirar excitação de algum lugar. No cérebro vivo ela existe: a cada neurônio chega o tempo todo uma atividade de fundo de milhares de outros. Suponha que o fundo mantenha o neurônio $Y$ ativo, e que a entrada $x$ o iniba por meio de um intermediário \nt{GABA}. Isso é um NÃO.

Mais útil, porém, é o \emph{veto}: ``$a$, mas não com $b$'':
\begin{empheq}[box=\keybox]{equation}
y \;=\; a \wedge \bar b .
\label{eq:veto}
\end{empheq}
O neurônio $Y$ recebe excitação de $a$ e inibição de um neurônio \nt{GABA} excitado por $b$. Aqui não é preciso fundo: a excitação vem do próprio $a$. Com vetos e OU monta-se tudo. Por exemplo, o OU exclusivo: $x\oplus y=(x\wedge\bar y)\vee(\bar x\wedge y)$, dois neurônios de veto, dois intermediários \nt{GABA} e um neurônio OU.

\begin{keyblock}
\begin{proposition}[Completude de \nt{GLUT}+\nt{GABA}]
\label{prop:gabacomplete}
Uma rede de neurônios \nt{GLUT} e \nt{GABA} pode calcular qualquer função lógica das entradas, e cada bit é transmitido por um único fio. Os sensores ``liga--desliga'' do capítulo~\ref{sec:glutcomputer} não são mais necessários para isso. Se a função deve dar um com todas as entradas caladas, é preciso uma fonte de atividade de fundo.
\end{proposition}
\end{keyblock}

Demonstração: o veto~\eqref{eq:veto} junto com o OU é funcionalmente completo se houver uma constante um, porque NÃO $x$ é o veto ``um, mas não com $x$''. E as funções que dão zero com todas as entradas caladas se montam com vetos e OU mesmo sem a constante um.

\section{Direção do movimento}

O veto no tempo, assim como o E no tempo no capítulo~\ref{sec:glutcomputer}, dá um detector de direção do movimento.

Sejam dois sensores vizinhos, $A$ e $B$, a uma distância $\Delta x$ um do outro. O neurônio de saída $Y$ recebe excitação de $B$ e inibição de $A$ por meio de um intermediário \nt{GABA}, com atraso $d$ (fig.~\ref{fig:direction}). Uma mancha de luz que corre com velocidade $v$ de $A$ para $B$ passa por $A$ no instante $0$, e a inibição chega a $Y$ no instante $d$; passa por $B$ no instante $\Delta x/v$, e nesse mesmo momento chega a excitação. Se esses instantes coincidem dentro da janela~\eqref{eq:window}, a inibição anula a excitação, e $Y$ fica calado. Isso acontece a uma velocidade de cerca de
\begin{empheq}[box=\keybox]{equation}
v^* \;=\; \frac{\Delta x}{d} .
\label{eq:vstar}
\end{empheq}
No movimento de $B$ para $A$, a excitação de $B$ chega no instante $0$, e a inibição de $A$ só no instante $\Delta x/v+d$, quando $Y$ já disparou.

\begin{figure}[t]
\centering
\begin{tikzpicture}[>=Stealth,
  nrn/.style={circle,draw,very thick,fill=blue!6,minimum size=0.9cm,inner sep=0pt},
  gab/.style={nrn,fill=red!10},
  inh/.style={-{Bar[width=7pt]},thick,red!70!black}]
\node[nrn] (A) at (0,2) {$A$};
\node[nrn] (B) at (3,2) {$B$};
\node[gab] (G) at (0.9,0.6) {\small \nt{GABA}};
\node[nrn] (Y) at (3,-0.6) {$Y$};
\draw[->,thick] (A) -- (G);
\draw[inh] (G) -- node[below left=-2pt] {\scriptsize atraso $d$} (Y);
\draw[->,thick] (B) -- (Y);
\draw[->,very thick,red!70!black] (Y) -- (4.6,-0.6) node[right,black] {\small saída};
\draw[->,thick,orange!85!black] (-0.6,2.9) -- (3.6,2.9) node[right,black] {\small calado};
\draw[<-,thick,orange!85!black] (-0.6,3.4) -- (3.6,3.4) node[right,black] {\small responde};
\foreach \yy/\dd in {2.9/-1,3.4/1}{
  \foreach \k in {1,2,3}{\draw[orange!70!yellow,line width=0.8pt] (1.5+\dd*0.12+\dd*0.13*\k,\yy+0.1-0.1*\k+0.1) -- ++(\dd*0.25,0);}
  \fill[yellow!70!orange,draw=orange!70!black] (1.5,\yy) circle (0.14);}
\node[font=\scriptsize,align=center] at (1.5,3.95) {mancha de luz};
\end{tikzpicture}
\caption{Detector de direção do movimento. $B$ excita $Y$, $A$ o inibe por meio de um intermediário \nt{GABA} com atraso $d$. No movimento de $A$ para $B$ com velocidade de cerca de $\Delta x/d$, a inibição anula a excitação; no movimento contrário, ela chega atrasada. A mancha amarela com tracinhos é a luz que corre sobre os sensores.}
\label{fig:direction}
\end{figure}

Na retina, a seletividade à direção do movimento parece ser criada exatamente assim: por uma inibição assimétrica vinda do lado de onde o movimento não deve provocar resposta~\cite{Barlow1965}.

\section{Subtração ou divisão}

Até agora, nossa inibição subtraía. Na verdade, ela é mais sutil.

Uma sinapse abre um canal, isto é, liga uma condutância. Numa sinapse excitatória, o potencial de equilíbrio fica muito acima do limiar, cerca de $70\mV$ acima do repouso: o canal aberto puxa a tensão para cima. Numa sinapse \nt{GABA} inibitória, o canal deixa passar cloreto, cujo potencial de equilíbrio fica perto do repouso, e o canal aberto puxa a tensão não para baixo, mas \emph{para o repouso}. Contando a tensão a partir do repouso, a tensão de regime de uma membrana com condutância de vazamento $g_L$, excitatória $g_E$ e inibitória $g_I$ é
\begin{empheq}[box=\keybox]{equation}
V_\infty \;=\; \frac{g_E\,E_E}{g_L+g_E+g_I} ,
\label{eq:shunt}
\end{empheq}
em que $E_E\approx70\mV$ é o potencial de equilíbrio da sinapse excitatória. É a lei de Ohm para um divisor de tensão: $g_E$ puxa para $E_E$, e $g_L$ e $g_I$, para zero.

$g_I$ não aparece no numerador com sinal de menos, e sim no denominador: a inibição não \emph{subtrai} da excitação, ela a \emph{divide} (fig.~\ref{fig:shunt}). Essa inibição se chama inibição por derivação: os canais de cloreto abertos curto-circuitam a membrana. Para a rede, é um controle de ganho. Se $g_I$ é proporcional à atividade total dos vizinhos, cada neurônio responde não à força absoluta de sua entrada, mas à sua fração na atividade total. Essa divisão pela atividade total aparece em muitas regiões do cérebro e é considerada uma de suas operações padrão~\cite{Carandini2012}: a mesma rede funciona tanto com entrada fraca quanto com entrada forte.

Uma ressalva: a fórmula~\eqref{eq:shunt} trata de um neurônio que não emite disparos. Num neurônio que dispara, a tensão fica o tempo todo perto do limiar, a corrente pela condutância inibitória é quase constante, e sobre a \emph{frequência} de disparos a derivação age principalmente por subtração~\cite{HoltKoch1997}. A frequência é dividida se ao neurônio forem acrescentadas ao mesmo tempo condutâncias de fundo excitatória e inibitória equilibradas, como no regime trêmulo e equilibrado do córtex vivo~\cite{Chance2002}. Assim, o cérebro obtém a divisão não de uma derivação isolada, mas da inibição junto com o ruído de fundo~\cite{Carandini2012}.

\begin{figure}[t]
\centering
\begin{tikzpicture}[>=Stealth,x=1.25cm,y=0.05cm]
\draw[->,gray!70!black] (0,0) -- (5.4,0) node[right,black] {\small $g_E/g_L$};
\draw[->,gray!70!black] (0,0) -- (0,68) node[above,black] {\small $V_\infty$, mV};
\foreach \v in {20,40,60}{\draw[gray!60!black] (-0.05,\v) -- (0.05,\v) node[left=3pt,font=\scriptsize] {\v};}
\foreach \x in {1,2,3,4,5}{\draw[gray!60!black] (\x,-1.5) -- (\x,1.5) node[below=5pt,font=\scriptsize] {\x};}
\draw[very thick,blue!60!black] plot[domain=0:5,samples=60] (\x,{70*\x/(1+\x)});
\draw[very thick,red!70!black] plot[domain=0:5,samples=60] (\x,{70*\x/(3+\x)});
\draw[thick,densely dashed,gray!50!black] plot[domain=0.56:5,samples=60] (\x,{70*\x/(1+\x)-25});
\node[right,font=\small,blue!60!black] at (5.02,58.3) {sem inibição};
\node[right,font=\small,red!70!black] at (5.02,45.5) {divide: $g_I=2g_L$};
\node[right,font=\small,gray!40!black] at (5.02,32.5) {subtrai $25\mV$};
\end{tikzpicture}
\caption{A inibição por derivação divide a tensão em vez de subtrair dela. A curva azul é a tensão de regime~\eqref{eq:shunt} sem inibição; a vermelha, com condutância inibitória $g_I=2g_L$. A vermelha é a mesma azul esticada três vezes na horizontal: como se a entrada fosse dividida por $1+g_I/g_L=3$. A tracejada é uma inibição que subtrai $25\mV$ constantes: a curva desce inteira, e a inclinação não muda.}
\label{fig:shunt}
\end{figure}

Há também uma segunda consequência. A constante de tempo da membrana é $\tau_{\text{ef}}=C/(g_L+g_E+g_I)$, e a inibição a encurta. Como a janela de coincidência~\eqref{eq:window} é proporcional a $\tau$, uma inibição que chega junto com a excitação \emph{estreita a janela}. O circuito em que a entrada excita o neurônio diretamente e, ao mesmo tempo, com atraso de um milissegundo, o inibe por meio de um intermediário \nt{GABA} é um dos mais comuns no cérebro: o neurônio só tem tempo de responder ao que chegou no primeiro ou segundo milissegundo~\cite{Pouille2001}.

\section{Escolha}

Agora o principal que faltava à rede \nt{GLUT}: escolher uma entre várias opções. Tomemos dois grupos de neurônios \nt{GLUT} com entradas $I_1$ e $I_2$. Cada um, por meio de seus intermediários \nt{GABA}, inibe o outro com força $\beta$ (fig.~\ref{fig:wta}). Para as frequências $r_1,r_2$ obtemos
\begin{equation}
\tau\,\frac{\dd r_1}{\dd t} = -r_1 + \bigl[I_1-\beta r_2\bigr]_+ ,\qquad
\tau\,\frac{\dd r_2}{\dd t} = -r_2 + \bigl[I_2-\beta r_1\bigr]_+ ,
\label{eq:wta}
\end{equation}
em que $[u]_+=\max(u,0)$: a frequência não pode ser negativa.

\begin{figure}[t]
\centering
\begin{tikzpicture}[>=Stealth,
  nrn/.style={circle,draw,very thick,fill=blue!6,minimum size=0.9cm,inner sep=0pt},
  gab/.style={nrn,fill=red!10,minimum size=0.75cm},
  inh/.style={-{Bar[width=7pt]},thick,red!70!black}]
\node[nrn] (E1) at (0,0) {$r_1$};
\node[nrn] (E2) at (4,0) {$r_2$};
\node[gab] (G1) at (1.4,1.1) {};
\node[gab] (G2) at (2.6,-1.1) {};
\draw[->,thick] (E1) -- (G1); \draw[inh] (G1) -- (E2);
\draw[->,thick] (E2) -- (G2); \draw[inh] (G2) -- (E1);
\draw[->,thick,blue!60!black] (-1.6,0) node[left,black] {$I_1$} -- (E1);
\draw[->,thick,blue!60!black] (5.6,0) node[right,black] {$I_2$} -- (E2);
\end{tikzpicture}
\caption{Escolha entre dois. Dois grupos de neurônios \nt{GLUT} (azuis) inibem um ao outro por meio de intermediários \nt{GABA} (vermelhos).}
\label{fig:wta}
\end{figure}

Enquanto os dois neurônios estão ativos, o sistema~\eqref{eq:wta} é linear, e seu comportamento é dado por dois autovalores: $-(1+\beta)/\tau$ para desvios iguais das duas frequências e $-(1-\beta)/\tau$ para desvios opostos. Com inibição fraca, $\beta<1$, os dois autovalores são negativos: os dois neurônios trabalham, e a inibição apenas abafa o mais fraco. Com inibição forte, $\beta>1$, o segundo autovalor fica positivo: qualquer diferença entre $r_1$ e $r_2$ cresce até que um neurônio se cale por completo.

\begin{keyblock}
\begin{proposition}[O vencedor leva tudo]
\label{prop:wta}
Se a inibição mútua é maior que um, $\beta>1$, o estado em que os dois grupos trabalham é instável. A rede chega a um estado em que um grupo trabalha e o outro fica calado. O vencedor, com frequência $r_1=I_1$, mantém a vitória enquanto $I_2<\beta I_1$.
\end{proposition}
\end{keyblock}

Verificamos isso num modelo. Com $\beta=0.5$ e entradas $1.0$ e $0.9$, os dois grupos trabalham, com frequências $0.73$ e $0.53$. Com $\beta=2$ e as mesmas entradas, o segundo grupo fica totalmente calado. Essa escolha tem memória: se depois as entradas forem trocadas, o vencedor anterior continua vencedor, porque $1.0<2\cdot0.9$.

Com cem grupos é igual: cada um inibe os demais, e sobrevive um.

\section{Reset da memória}

A memória do capítulo~\ref{sec:glutcomputer} só se desligava pela fadiga. Vamos acrescentar a ela um neurônio \nt{GABA} que é excitado por um sinal de ``reset'' e inibe o grupo. A inibição abaixa a curva $f(wr+I)$ da fig.~\ref{fig:bistable}, a interseção de cima some e o grupo se apaga, e fica no estado de baixo quando o reset termina. Agora a memória é gravada por um sinal e apagada por outro, como num flip-flop com entradas de set e reset.

Mais elegante ainda é ligar dois desses grupos por inibição mútua, como na fig.~\ref{fig:wta}, e dar a cada um realimentação sobre si mesmo. Um sinal na entrada de um dos grupos leva a célula ao estado dele, e a célula lembra a última decisão. É o análogo direto de um latch de dois elementos NOU acoplados em cruz.

\section{Estabilidade e ritmo}

Resta o terceiro mal da rede \nt{GLUT}: a avalanche. Tomemos um grupo de neurônios \nt{GLUT} com realimentação sobre si mesmo $w_{EE}$ e um grupo de neurônios \nt{GABA} que o primeiro excita com força $w_{IE}$ e que inibem o primeiro com força $w_{EI}$. Para as frequências $E$ e $I$ obtemos
\begin{equation}
\tau_E\frac{\dd E}{\dd t} = -E + w_{EE}E - w_{EI}I + h, \qquad
\tau_I\frac{\dd I}{\dd t} = -I + w_{IE}E ,
\label{eq:ei}
\end{equation}
em que $h$ é a entrada externa~\cite{WilsonCowan1972}. Sem inibição, com $w_{EE}>1$, a atividade cresce sem limite: cada disparo gera mais de um disparo seguinte. Com inibição, a frequência de regime
\begin{equation}
E^* \;=\; \frac{h}{1-w_{EE}+w_{EI}w_{IE}}
\label{eq:eistar}
\end{equation}
é finita se o denominador for positivo. Mas isso não basta: o equilíbrio ainda precisa atrair. Da análise linear de~\eqref{eq:ei} saem duas condições.

\begin{keyblock}
\begin{proposition}[Condições de estabilidade]
\label{prop:ei}
O equilíbrio~\eqref{eq:eistar} é estável se a inibição for
\begin{enumerate}
\item[(i)] \emph{forte o bastante}: $w_{EI}\,w_{IE} > w_{EE}-1$;
\item[(ii)] \emph{rápida o bastante}: $\tau_I < \tau_E/(w_{EE}-1)$.
\end{enumerate}
Se (i) é violada, a atividade cresce em avalanche. Se (i) vale mas (ii) é violada, a inibição chega atrasada, e a atividade começa a oscilar.
\end{proposition}
\end{keyblock}

A condição~(i) é o determinante da matriz do sistema~\eqref{eq:ei}; a condição~(ii), o seu traço. O sentido da segunda é claro para quem já ajustou um controlador: uma realimentação atrasada não amortece o desvio, mas o amplifica.

\begin{figure}[t]
\centering
\begin{tikzpicture}[>=Stealth]
\draw[->,gray!70!black] (0,0) -- (10.9,0) node[below left,black] {\small $t$, ms};
\draw[->,gray!70!black] (0,0) -- (0,3.3) node[left,black] {\small $E$};
\foreach \t in {0,50,100,150}{\draw[gray!70!black] (\t*0.07,0.06) -- (\t*0.07,-0.06) node[below,font=\scriptsize] {\t};}
\draw[very thick,gray!60!black,densely dashed] plot coordinates {(0.000,0.000) (0.035,0.071) (0.070,0.144) (0.105,0.218) (0.140,0.294) (0.175,0.373) (0.210,0.453) (0.245,0.535) (0.280,0.620) (0.315,0.706) (0.350,0.795) (0.385,0.886) (0.420,0.979) (0.455,1.075) (0.490,1.173) (0.525,1.274) (0.560,1.377) (0.595,1.482) (0.630,1.591) (0.665,1.702) (0.700,1.816) (0.735,1.933) (0.770,2.052) (0.805,2.175) (0.840,2.301) (0.875,2.430) (0.910,2.563) (0.945,2.698) (0.980,2.838) (1.015,2.980) (1.050,3.080) (1.085,3.080) (1.120,3.080) (1.155,3.080) (1.190,3.080) (1.225,3.080) (1.260,3.080) (1.295,3.080) (1.330,3.080) (1.365,3.080) (1.400,3.080) (1.435,3.080) (1.470,3.080) (1.505,3.080) (1.540,3.080) (1.575,3.080) (1.610,3.080) (1.645,3.080) (1.680,3.080) (1.715,3.080) (1.750,3.080) (1.785,3.080) (1.820,3.080) (1.855,3.080) (1.890,3.080) (1.925,3.080) (1.960,3.080) (1.995,3.080) (2.030,3.080) (2.065,3.080) (2.100,3.080) (2.135,3.080) (2.170,3.080) (2.205,3.080) (2.240,3.080) (2.275,3.080) (2.310,3.080) (2.345,3.080) (2.380,3.080) (2.415,3.080) (2.450,3.080) (2.485,3.080) (2.520,3.080) (2.555,3.080) (2.590,3.080) (2.625,3.080) (2.660,3.080) (2.695,3.080) (2.730,3.080) (2.765,3.080) (2.800,3.080) (2.835,3.080) (2.870,3.080) (2.905,3.080) (2.940,3.080) (2.975,3.080) (3.010,3.080) (3.045,3.080) (3.080,3.080) (3.115,3.080) (3.150,3.080) (3.185,3.080) (3.220,3.080) (3.255,3.080) (3.290,3.080) (3.325,3.080) (3.360,3.080) (3.395,3.080) (3.430,3.080) (3.465,3.080) (3.500,3.080) (3.535,3.080) (3.570,3.080) (3.605,3.080) (3.640,3.080) (3.675,3.080) (3.710,3.080) (3.745,3.080) (3.780,3.080) (3.815,3.080) (3.850,3.080) (3.885,3.080) (3.920,3.080) (3.955,3.080) (3.990,3.080) (4.025,3.080) (4.060,3.080) (4.095,3.080) (4.130,3.080) (4.165,3.080) (4.200,3.080) (4.235,3.080) (4.270,3.080) (4.305,3.080) (4.340,3.080) (4.375,3.080) (4.410,3.080) (4.445,3.080) (4.480,3.080) (4.515,3.080) (4.550,3.080) (4.585,3.080) (4.620,3.080) (4.655,3.080) (4.690,3.080) (4.725,3.080) (4.760,3.080) (4.795,3.080) (4.830,3.080) (4.865,3.080) (4.900,3.080) (4.935,3.080) (4.970,3.080) (5.005,3.080) (5.040,3.080) (5.075,3.080) (5.110,3.080) (5.145,3.080) (5.180,3.080) (5.215,3.080) (5.250,3.080) (5.285,3.080) (5.320,3.080) (5.355,3.080) (5.390,3.080) (5.425,3.080) (5.460,3.080) (5.495,3.080) (5.530,3.080) (5.565,3.080) (5.600,3.080) (5.635,3.080) (5.670,3.080) (5.705,3.080) (5.740,3.080) (5.775,3.080) (5.810,3.080) (5.845,3.080) (5.880,3.080) (5.915,3.080) (5.950,3.080) (5.985,3.080) (6.020,3.080) (6.055,3.080) (6.090,3.080) (6.125,3.080) (6.160,3.080) (6.195,3.080) (6.230,3.080) (6.265,3.080) (6.300,3.080) (6.335,3.080) (6.370,3.080) (6.405,3.080) (6.440,3.080) (6.475,3.080) (6.510,3.080) (6.545,3.080) (6.580,3.080) (6.615,3.080) (6.650,3.080) (6.685,3.080) (6.720,3.080) (6.755,3.080) (6.790,3.080) (6.825,3.080) (6.860,3.080) (6.895,3.080) (6.930,3.080) (6.965,3.080) (7.000,3.080) (7.035,3.080) (7.070,3.080) (7.105,3.080) (7.140,3.080) (7.175,3.080) (7.210,3.080) (7.245,3.080) (7.280,3.080) (7.315,3.080) (7.350,3.080) (7.385,3.080) (7.420,3.080) (7.455,3.080) (7.490,3.080) (7.525,3.080) (7.560,3.080) (7.595,3.080) (7.630,3.080) (7.665,3.080) (7.700,3.080) (7.735,3.080) (7.770,3.080) (7.805,3.080) (7.840,3.080) (7.875,3.080) (7.910,3.080) (7.945,3.080) (7.980,3.080) (8.015,3.080) (8.050,3.080) (8.085,3.080) (8.120,3.080) (8.155,3.080) (8.190,3.080) (8.225,3.080) (8.260,3.080) (8.295,3.080) (8.330,3.080) (8.365,3.080) (8.400,3.080) (8.435,3.080) (8.470,3.080) (8.505,3.080) (8.540,3.080) (8.575,3.080) (8.610,3.080) (8.645,3.080) (8.680,3.080) (8.715,3.080) (8.750,3.080) (8.785,3.080) (8.820,3.080) (8.855,3.080) (8.890,3.080) (8.925,3.080) (8.960,3.080) (8.995,3.080) (9.030,3.080) (9.065,3.080) (9.100,3.080) (9.135,3.080) (9.170,3.080) (9.205,3.080) (9.240,3.080) (9.275,3.080) (9.310,3.080) (9.345,3.080) (9.380,3.080) (9.415,3.080) (9.450,3.080) (9.485,3.080) (9.520,3.080) (9.555,3.080) (9.590,3.080) (9.625,3.080) (9.660,3.080) (9.695,3.080) (9.730,3.080) (9.765,3.080) (9.800,3.080) (9.835,3.080) (9.870,3.080) (9.905,3.080) (9.940,3.080) (9.975,3.080) (10.010,3.080) (10.045,3.080) (10.080,3.080) (10.115,3.080) (10.150,3.080) (10.185,3.080) (10.220,3.080) (10.255,3.080) (10.290,3.080) (10.325,3.080) (10.360,3.080) (10.395,3.080) (10.430,3.080) (10.465,3.080) (10.500,3.080)};
\draw[very thick,blue!60!black] plot coordinates {(0.000,0.000) (0.035,0.071) (0.070,0.141) (0.105,0.211) (0.140,0.279) (0.175,0.344) (0.210,0.406) (0.245,0.465) (0.280,0.519) (0.315,0.570) (0.350,0.616) (0.385,0.658) (0.420,0.697) (0.455,0.731) (0.490,0.763) (0.525,0.790) (0.560,0.815) (0.595,0.836) (0.630,0.855) (0.665,0.871) (0.700,0.886) (0.735,0.898) (0.770,0.908) (0.805,0.916) (0.840,0.924) (0.875,0.929) (0.910,0.934) (0.945,0.938) (0.980,0.941) (1.015,0.943) (1.050,0.945) (1.085,0.946) (1.120,0.946) (1.155,0.947) (1.190,0.947) (1.225,0.947) (1.260,0.946) (1.295,0.946) (1.330,0.945) (1.365,0.944) (1.400,0.944) (1.435,0.943) (1.470,0.942) (1.505,0.941) (1.540,0.941) (1.575,0.940) (1.610,0.939) (1.645,0.939) (1.680,0.938) (1.715,0.937) (1.750,0.937) (1.785,0.936) (1.820,0.936) (1.855,0.936) (1.890,0.935) (1.925,0.935) (1.960,0.935) (1.995,0.934) (2.030,0.934) (2.065,0.934) (2.100,0.934) (2.135,0.934) (2.170,0.934) (2.205,0.934) (2.240,0.933) (2.275,0.933) (2.310,0.933) (2.345,0.933) (2.380,0.933) (2.415,0.933) (2.450,0.933) (2.485,0.933) (2.520,0.933) (2.555,0.933) (2.590,0.933) (2.625,0.933) (2.660,0.933) (2.695,0.933) (2.730,0.933) (2.765,0.933) (2.800,0.933) (2.835,0.933) (2.870,0.933) (2.905,0.933) (2.940,0.933) (2.975,0.933) (3.010,0.933) (3.045,0.933) (3.080,0.933) (3.115,0.933) (3.150,0.933) (3.185,0.933) (3.220,0.933) (3.255,0.933) (3.290,0.933) (3.325,0.933) (3.360,0.933) (3.395,0.933) (3.430,0.933) (3.465,0.933) (3.500,0.933) (3.535,0.933) (3.570,0.933) (3.605,0.933) (3.640,0.933) (3.675,0.933) (3.710,0.933) (3.745,0.933) (3.780,0.933) (3.815,0.933) (3.850,0.933) (3.885,0.933) (3.920,0.933) (3.955,0.933) (3.990,0.933) (4.025,0.933) (4.060,0.933) (4.095,0.933) (4.130,0.933) (4.165,0.933) (4.200,0.933) (4.235,0.933) (4.270,0.933) (4.305,0.933) (4.340,0.933) (4.375,0.933) (4.410,0.933) (4.445,0.933) (4.480,0.933) (4.515,0.933) (4.550,0.933) (4.585,0.933) (4.620,0.933) (4.655,0.933) (4.690,0.933) (4.725,0.933) (4.760,0.933) (4.795,0.933) (4.830,0.933) (4.865,0.933) (4.900,0.933) (4.935,0.933) (4.970,0.933) (5.005,0.933) (5.040,0.933) (5.075,0.933) (5.110,0.933) (5.145,0.933) (5.180,0.933) (5.215,0.933) (5.250,0.933) (5.285,0.933) (5.320,0.933) (5.355,0.933) (5.390,0.933) (5.425,0.933) (5.460,0.933) (5.495,0.933) (5.530,0.933) (5.565,0.933) (5.600,0.933) (5.635,0.933) (5.670,0.933) (5.705,0.933) (5.740,0.933) (5.775,0.933) (5.810,0.933) (5.845,0.933) (5.880,0.933) (5.915,0.933) (5.950,0.933) (5.985,0.933) (6.020,0.933) (6.055,0.933) (6.090,0.933) (6.125,0.933) (6.160,0.933) (6.195,0.933) (6.230,0.933) (6.265,0.933) (6.300,0.933) (6.335,0.933) (6.370,0.933) (6.405,0.933) (6.440,0.933) (6.475,0.933) (6.510,0.933) (6.545,0.933) (6.580,0.933) (6.615,0.933) (6.650,0.933) (6.685,0.933) (6.720,0.933) (6.755,0.933) (6.790,0.933) (6.825,0.933) (6.860,0.933) (6.895,0.933) (6.930,0.933) (6.965,0.933) (7.000,0.933) (7.035,0.933) (7.070,0.933) (7.105,0.933) (7.140,0.933) (7.175,0.933) (7.210,0.933) (7.245,0.933) (7.280,0.933) (7.315,0.933) (7.350,0.933) (7.385,0.933) (7.420,0.933) (7.455,0.933) (7.490,0.933) (7.525,0.933) (7.560,0.933) (7.595,0.933) (7.630,0.933) (7.665,0.933) (7.700,0.933) (7.735,0.933) (7.770,0.933) (7.805,0.933) (7.840,0.933) (7.875,0.933) (7.910,0.933) (7.945,0.933) (7.980,0.933) (8.015,0.933) (8.050,0.933) (8.085,0.933) (8.120,0.933) (8.155,0.933) (8.190,0.933) (8.225,0.933) (8.260,0.933) (8.295,0.933) (8.330,0.933) (8.365,0.933) (8.400,0.933) (8.435,0.933) (8.470,0.933) (8.505,0.933) (8.540,0.933) (8.575,0.933) (8.610,0.933) (8.645,0.933) (8.680,0.933) (8.715,0.933) (8.750,0.933) (8.785,0.933) (8.820,0.933) (8.855,0.933) (8.890,0.933) (8.925,0.933) (8.960,0.933) (8.995,0.933) (9.030,0.933) (9.065,0.933) (9.100,0.933) (9.135,0.933) (9.170,0.933) (9.205,0.933) (9.240,0.933) (9.275,0.933) (9.310,0.933) (9.345,0.933) (9.380,0.933) (9.415,0.933) (9.450,0.933) (9.485,0.933) (9.520,0.933) (9.555,0.933) (9.590,0.933) (9.625,0.933) (9.660,0.933) (9.695,0.933) (9.730,0.933) (9.765,0.933) (9.800,0.933) (9.835,0.933) (9.870,0.933) (9.905,0.933) (9.940,0.933) (9.975,0.933) (10.010,0.933) (10.045,0.933) (10.080,0.933) (10.115,0.933) (10.150,0.933) (10.185,0.933) (10.220,0.933) (10.255,0.933) (10.290,0.933) (10.325,0.933) (10.360,0.933) (10.395,0.933) (10.430,0.933) (10.465,0.933) (10.500,0.933)};
\draw[very thick,red!70!black] plot coordinates {(0.000,0.000) (0.035,0.071) (0.070,0.143) (0.105,0.217) (0.140,0.293) (0.175,0.370) (0.210,0.449) (0.245,0.528) (0.280,0.609) (0.315,0.691) (0.350,0.774) (0.385,0.858) (0.420,0.943) (0.455,1.028) (0.490,1.114) (0.525,1.200) (0.560,1.287) (0.595,1.374) (0.630,1.462) (0.665,1.549) (0.700,1.636) (0.735,1.724) (0.770,1.810) (0.805,1.897) (0.840,1.983) (0.875,2.068) (0.910,2.153) (0.945,2.237) (0.980,2.320) (1.015,2.402) (1.050,2.483) (1.085,2.563) (1.120,2.641) (1.155,2.717) (1.190,2.792) (1.225,2.866) (1.260,2.937) (1.295,3.007) (1.330,3.075) (1.365,3.080) (1.400,3.080) (1.435,3.080) (1.470,3.080) (1.505,3.080) (1.540,3.080) (1.575,3.080) (1.610,3.080) (1.645,3.080) (1.680,3.080) (1.715,3.080) (1.750,3.080) (1.785,3.080) (1.820,3.080) (1.855,3.080) (1.890,3.080) (1.925,3.080) (1.960,3.080) (1.995,3.080) (2.030,3.080) (2.065,3.080) (2.100,3.080) (2.135,3.080) (2.170,3.080) (2.205,3.080) (2.240,3.080) (2.275,3.080) (2.310,3.080) (2.345,3.080) (2.380,3.080) (2.415,3.080) (2.450,3.080) (2.485,3.080) (2.520,3.080) (2.555,3.080) (2.590,3.080) (2.625,3.080) (2.660,3.080) (2.695,3.080) (2.730,3.080) (2.765,3.080) (2.800,3.080) (2.835,3.052) (2.870,2.973) (2.905,2.892) (2.940,2.807) (2.975,2.719) (3.010,2.629) (3.045,2.535) (3.080,2.438) (3.115,2.339) (3.150,2.238) (3.185,2.133) (3.220,2.029) (3.255,1.930) (3.290,1.836) (3.325,1.747) (3.360,1.661) (3.395,1.580) (3.430,1.503) (3.465,1.430) (3.500,1.360) (3.535,1.294) (3.570,1.231) (3.605,1.171) (3.640,1.113) (3.675,1.059) (3.710,1.007) (3.745,0.958) (3.780,0.911) (3.815,0.867) (3.850,0.825) (3.885,0.784) (3.920,0.746) (3.955,0.710) (3.990,0.675) (4.025,0.642) (4.060,0.611) (4.095,0.581) (4.130,0.553) (4.165,0.526) (4.200,0.500) (4.235,0.476) (4.270,0.452) (4.305,0.430) (4.340,0.409) (4.375,0.389) (4.410,0.370) (4.445,0.352) (4.480,0.335) (4.515,0.319) (4.550,0.303) (4.585,0.288) (4.620,0.274) (4.655,0.261) (4.690,0.248) (4.725,0.236) (4.760,0.225) (4.795,0.214) (4.830,0.203) (4.865,0.193) (4.900,0.184) (4.935,0.175) (4.970,0.166) (5.005,0.158) (5.040,0.151) (5.075,0.143) (5.110,0.136) (5.145,0.130) (5.180,0.123) (5.215,0.117) (5.250,0.111) (5.285,0.106) (5.320,0.101) (5.355,0.096) (5.390,0.091) (5.425,0.087) (5.460,0.083) (5.495,0.079) (5.530,0.075) (5.565,0.071) (5.600,0.068) (5.635,0.064) (5.670,0.061) (5.705,0.058) (5.740,0.055) (5.775,0.053) (5.810,0.050) (5.845,0.048) (5.880,0.045) (5.915,0.043) (5.950,0.041) (5.985,0.039) (6.020,0.037) (6.055,0.035) (6.090,0.034) (6.125,0.032) (6.160,0.030) (6.195,0.029) (6.230,0.027) (6.265,0.026) (6.300,0.025) (6.335,0.024) (6.370,0.022) (6.405,0.021) (6.440,0.021) (6.475,0.022) (6.510,0.024) (6.545,0.027) (6.580,0.032) (6.615,0.037) (6.650,0.044) (6.685,0.052) (6.720,0.061) (6.755,0.071) (6.790,0.083) (6.825,0.095) (6.860,0.109) (6.895,0.124) (6.930,0.140) (6.965,0.157) (7.000,0.176) (7.035,0.195) (7.070,0.215) (7.105,0.237) (7.140,0.260) (7.175,0.283) (7.210,0.308) (7.245,0.333) (7.280,0.360) (7.315,0.388) (7.350,0.416) (7.385,0.445) (7.420,0.476) (7.455,0.507) (7.490,0.539) (7.525,0.571) (7.560,0.604) (7.595,0.638) (7.630,0.673) (7.665,0.708) (7.700,0.744) (7.735,0.781) (7.770,0.818) (7.805,0.855) (7.840,0.893) (7.875,0.931) (7.910,0.969) (7.945,1.008) (7.980,1.047) (8.015,1.086) (8.050,1.126) (8.085,1.165) (8.120,1.204) (8.155,1.244) (8.190,1.283) (8.225,1.322) (8.260,1.362) (8.295,1.400) (8.330,1.439) (8.365,1.477) (8.400,1.515) (8.435,1.553) (8.470,1.590) (8.505,1.626) (8.540,1.662) (8.575,1.698) (8.610,1.733) (8.645,1.767) (8.680,1.800) (8.715,1.832) (8.750,1.864) (8.785,1.895) (8.820,1.924) (8.855,1.953) (8.890,1.981) (8.925,2.007) (8.960,2.033) (8.995,2.057) (9.030,2.080) (9.065,2.102) (9.100,2.123) (9.135,2.142) (9.170,2.160) (9.205,2.177) (9.240,2.192) (9.275,2.205) (9.310,2.217) (9.345,2.228) (9.380,2.237) (9.415,2.245) (9.450,2.251) (9.485,2.255) (9.520,2.258) (9.555,2.259) (9.590,2.259) (9.625,2.256) (9.660,2.253) (9.695,2.247) (9.730,2.240) (9.765,2.231) (9.800,2.220) (9.835,2.208) (9.870,2.194) (9.905,2.178) (9.940,2.160) (9.975,2.141) (10.010,2.120) (10.045,2.098) (10.080,2.074) (10.115,2.048) (10.150,2.021) (10.185,1.992) (10.220,1.961) (10.255,1.929) (10.290,1.895) (10.325,1.860) (10.360,1.824) (10.395,1.786) (10.430,1.746) (10.465,1.706) (10.500,1.663)};

\node[font=\small,gray!40!black,anchor=west] at (0.9,3.15) {sem inibição: avalanche};
\node[font=\small,blue!60!black,anchor=west] at (7.4,1.25) {inibição rápida};
\node[font=\small,red!70!black,anchor=west] at (7.4,2.55) {inibição lenta};
\end{tikzpicture}
\caption{Frequência de um grupo \nt{GLUT} com realimentação $w_{EE}=1.5$, $\tau_E=10\ms$, após a entrada ser ligada. Sem inibição (tracejado), a atividade cresce sem limite. Com inibição de força $w_{EI}w_{IE}=2$ e $\tau_I=2\ms$ (curva azul), ela se estabiliza no nível $E^*=h/(1-1.5+2)=0.67h$. Com a mesma inibição, mas lenta, $\tau_I=30\ms$ (curva vermelha), a atividade oscila.}
\label{fig:ei}
\end{figure}

Considera-se que o córtex trabalha exatamente nesse regime: sua própria realimentação é forte o bastante para cair em avalanche sem inibição, e só a inibição rápida o mantém no ponto de operação~\cite{Tsodyks1997isn}.

Esse regime tem uma assinatura paradoxal pela qual pode ser reconhecido de fora~\cite{Tsodyks1997isn}. Acrescentemos em~\eqref{eq:ei} uma entrada externa $h_I$ diretamente ao grupo \nt{GABA}. As frequências de regime passam a ser
\begin{equation}
E^* \;=\; \frac{h-w_{EI}\,h_I}{D},\qquad
I^* \;=\; \frac{w_{IE}\,h+(1-w_{EE})\,h_I}{D},\qquad
D \;=\; 1-w_{EE}+w_{EI}w_{IE} .
\label{eq:isn}
\end{equation}
Se $w_{EE}>1$, o coeficiente de $h_I$ na segunda fórmula é negativo: empurre os neurônios \nt{GABA}, e sua frequência \emph{cai}. A causa está no laço. A inibição extra abafa o grupo \nt{GLUT}, este excita menos os neurônios \nt{GABA}, e eles perdem mais do que ganharam. Com os números da fig.~\ref{fig:ei} ($w_{EE}=1.5$, $w_{EI}w_{IE}=2$, $D=1.5$), cada unidade de entrada acrescentada reduz $I^*$ em um terço de unidade. Essa assinatura foi encontrada no córtex: quando a resposta de neurônios do córtex visual é suprimida por um estímulo ao redor, a corrente inibitória que eles recebem não cresce, mas cai junto com a excitatória~\cite{Ozeki2009}. Mais tarde a mesma assinatura foi encontrada em diferentes áreas e camadas do córtex~\cite{Sanzeni2020}.

Oscilações quando a condição~(ii) é violada também ocorrem no cérebro. O laço ``\nt{GLUT} excita \nt{GABA}, \nt{GABA} inibe \nt{GLUT}'' com atraso de alguns milissegundos gera um ritmo com período da ordem de $12$--$50\ms$ (frequência de $20$--$80$\,Hz), e esses ritmos rápidos no córtex são bem conhecidos~\cite{Cardin2009,Buzsaki2012}. Não é o clock de um computador: o ritmo é local e só surge quando a rede está ativa. Mas, ao longo de alguns ciclos, ele divide o tempo em janelas nas quais os neurônios podem disparar.

\section{Resumo}

\begin{table}[t]
\centering
\small
\begin{tabular}{>{\raggedright}p{3.3cm}>{\raggedright}p{4.4cm}>{\raggedright\arraybackslash}p{4.4cm}}
\toprule
& só \nt{GLUT} & \nt{GLUT} + \nt{GABA} \\
\midrule
NÃO & só por sensores ``liga--desliga'' & veto $a\wedge\bar b$; NÃO com atividade de fundo \\
fios por bit & dois & um \\
direção do movimento & não & veto com atraso \\
controle de ganho & não & divisão: derivação junto com ruído de fundo \\
janela de coincidência & $\sim\tau$ & estreitada pela inibição \\
escolha de um entre muitos & não & inibição mútua, $\beta>1$ \\
reset da memória & só por fadiga & por sinal via \nt{GABA} \\
estabilidade & só por fadiga & realimentação negativa rápida \\
ritmo & desnecessário & surge com inibição lenta \\
\bottomrule
\end{tabular}
\caption{O que o \nt{GABA} acrescenta a uma rede de neurônios \nt{GLUT}.}
\label{tab:gaba}
\end{table}

O \nt{GABA} quase não acrescenta à rede novos \emph{cálculos} (tudo isso também podia ser calculado com sensores de dois fios), mas acrescenta \emph{controle} (tabela~\ref{tab:gaba}): escolha, reset, controle de ganho e estabilidade. No cérebro, a excitação leva os dados, e a inibição decide quais dados passam, quando e com que volume. Para um em cada quatro ou cinco neurônios do córtex, é uma divisão de trabalho muito sensata.

\section{Exercícios}

\begin{exercise}[\lvlA]
\label{ex:03:1}
\emph{A derivação em números.} $E_E=70\mV$. Por~\eqref{eq:shunt}, encontre $V_\infty$ para $g_E=g_L$ e $4g_L$, sem inibição e com derivação $g_I=2g_L$. Que $V_\infty$ em $g_E=4g_L$ daria uma subtração que, em $g_E=g_L$, retira tanto quanto a derivação? Quantas vezes a derivação estreita a janela de coincidência?
\end{exercise}

\begin{exercise}[\lvlB]
\label{ex:03:2}
\emph{Dedução das condições de estabilidade.} Encontre o traço e o determinante da matriz do sistema linear~\eqref{eq:ei} e deduza deles as duas condições da proposição~\ref{prop:ei}. Na fronteira da condição~(ii), o equilíbrio perde estabilidade de forma oscilatória. Encontre o período dessas oscilações com os números da fig.~\ref{fig:ei}.
\end{exercise}

\begin{exercise}[\lvlB]
\label{ex:03:3}
\emph{Ritmo a partir do atraso.} Troquemos a inibição lenta do modelo~\eqref{eq:ei} por uma rápida com atraso $d$: $\tau_E\dot E=-E(t)-GE(t-d)+h$. Com $\tau_E=10\ms$, encontre o menor atraso $d_c$ que produz oscilações e seu período para $G=2$ e $G=5$. Eles caem nos ritmos rápidos do córtex, $12$--$50\ms$, ao contrário do exercício~\ref{ex:03:2}?
\end{exercise}

\begin{exercise}[\lvlB]
\label{ex:03:4}
\emph{Modelagem: escolha com memória.} Integre~\eqref{eq:wta} com $\beta=2$, $\tau=1$. (a)~Com $I_1=1$, suba $I_2$ lentamente de $0$ a $3$ e desça de volta: em que valores de $I_2$ a rede comuta? (b)~Quanto cresce o tempo de decisão a partir do silêncio quando a diferença das entradas $\varepsilon$ diminui dez vezes?
\end{exercise}

\begin{exercise}[\lvlB]
\label{ex:03:5}
\emph{Projeto: detector para uma faixa de velocidades.} O detector da fig.~\ref{fig:direction} deve ficar calado no movimento de $A$ para $B$ com tempo de percurso de $5$--$20\ms$. Um intermediário \nt{GABA} com atraso $d_k$ veta $Y$ no intervalo $[d_k,\,d_k+5\ms]$ após o disparo de $A$; a excitação de $B$ chega de imediato. Quantos intermediários são necessários, e com que atrasos?
\end{exercise}

\begin{exercise}[\lvlB]
\label{ex:03:6}
\emph{Vetos e fundo.} Quantos neurônios exige a paridade $x\oplus y\oplus z$ no esquema geral ``intermediário \nt{GABA} por entrada, veto para cada combinação com $f=1$, OU comum''? Construa-a com dois neurônios, \nt{GABA} e \nt{GLUT}, que recebem as entradas diretamente, indicando pesos e limiares.
\end{exercise}

\begin{exercise}[\lvlC]
\label{ex:03:7}
\emph{Projeto: escolha entre cem.} Cada um de $100$ grupos \nt{GLUT} deve inibir igualmente todos os outros, mas não a si mesmo. Intermediários \nt{GABA} lineares são excitados pelos grupos e os inibem; os grupos podem excitar uns aos outros, mas não a si mesmos. Qual é o menor número de intermediários? E se não houver nenhuma conexão excitatória direta?
\end{exercise}

\begin{exercise}[\lvlC]
\label{ex:03:8}
\emph{Pesquisa: quando a inibição é paradoxal.} No modelo da fig.~\ref{fig:ei} ($w_{EI}=2$, $w_{IE}=1$), o grupo \nt{GLUT} recebeu a função de transferência $f(u)=[u]_+^2$, e o grupo \nt{GABA} continuou linear. Acima de que entrada $h_c$ um acréscimo ao grupo \nt{GABA} reduz a própria frequência dele? Verifique por simulação.
\end{exercise}
